% TOP_PROBLEM: {"id": 2100405, "title": "Open Problems in Integrable Systems — Geometry of caustics, invariant surfaces, and commuting billiard maps", "category_id": 11, "category": {"id": 11, "name": "dynamical_systems", "display_name": "Dynamical Systems", "description": "Problems about long-term behavior of deterministic systems, Hamiltonian dynamics, and periodic orbits.", "slug": "dynamical-systems", "order_index": 11, "created_at": "2026-07-31T15:26:25.671Z"}, "status": "open"}
% FIRST_POSTED: null
% ENTRY_KIND: statement_only
% Imported from: batch08.tex; UnsolvedMath ID 2100405
% Compile twice: lualatex 2100405.tex
\documentclass[11pt]{article}
\usepackage{fontspec}
\setmainfont{Latin Modern Roman}
\usepackage{amsmath,amssymb,mathtools,unicode-math}
\setmathfont{Latin Modern Math}
\usepackage{xurl,hyperref}
\hypersetup{hidelinks,pdftitle={UnsolvedMath Problem 2100405}}
\newcommand{\sref}[2]{\hyperlink{source:#1:#2}{[S#2]}}
\newcommand{\cataloguescope}[1]{\paragraph{Mathematical question.}#1}
\begin{document}
% UnsolvedMath ID 2100405; source catalogue code AMR-020-0405
\setcounter{section}{2100404}
\section{Open Problems in Integrable Systems — Geometry of caustics, invariant surfaces, and commuting billiard maps}\label{2100405}
{\small\sffamily UnsolvedMath ID 2100405 · Dynamical Systems}
\subsection{Definitions and mathematical statement}

\paragraph{Shared notation.}$\mathbb N=\{1,2,\ldots\}$, and $\mathbb Z,\mathbb Q,\mathbb R,\mathbb C$ denote the integers, rationals, reals and complex numbers. Any other conventions are specified locally.


For a bounded smooth strictly convex planar domain $\Omega$, the inner billiard reflects unit-speed straight trajectories at $\partial\Omega$, preserving the tangential velocity and reversing its normal component. A caustic is an envelope of an invariant one-parameter family of rays: tangency to it is preserved by reflection. A $p/q$-rational caustic corresponds to an invariant circle of the billiard map consisting of periodic points of rotation number $p/q$. The outer billiard map sends a point outside a smooth convex curve to its reflection through the chosen tangency point, using one consistent orientation of the tangent. For the outer map, rational caustics likewise mean invariant circles consisting of periodic points of a prescribed rational rotation number.
\paragraph{Statement recorded in the supplied source.}
Does there exist a planar billiard table that is not an ellipse and has rational caustics with two distinct rotation numbers? Ask the same existence question for planar outer billiards. The source includes the $1/2$ back-and-forth invariant family in its rational-caustic convention for inner billiards. \sref{2100405}{2}
\subsection{Short English statement}
Can a nonelliptic inner or outer billiard possess two rational caustics with different rotation numbers?
\subsection{Sources}
\begin{description}
\item[\hypertarget{source:2100405:1}{S1}] ULAM.AI and the UnsolvedMath Contributors, UnsolvedMath: A Curated Collection of Open Mathematics Problems, record 2100405 (AMR-020-0405). CC BY 4.0. \url{https://huggingface.co/datasets/ulamai/UnsolvedMath}.
\item[\hypertarget{source:2100405:2}{S2}] A. Bolsinov, V. S. Matveev, E. Miranda and S. Tabachnikov, Open Problems, Questions, and Challenges in Finite-Dimensional Integrable Systems (2018), arXiv:1804.03737, Question 4.7. \url{https://arxiv.org/abs/1804.03737}.
\end{description}
\label{end:2100405}
\end{document}
